\documentclass[11pt]{article}
\usepackage[letterpaper,margin=1in]{geometry}
\usepackage[T1]{fontenc}
\usepackage[utf8]{inputenc}
\usepackage{amsmath,amssymb,amsthm,lmodern,microtype,booktabs}
\usepackage[hidelinks]{hyperref}

\newtheorem{theorem}{Theorem}
\newtheorem{lemma}[theorem]{Lemma}
\newtheorem{proposition}[theorem]{Proposition}

\title{Berge--Fulkerson Covers Across a Hexagon}
\author{}
\date{}

\begin{document}
\maketitle

\begin{abstract}
Let $G$ be a cubic graph containing an induced $6$-cycle (a hexagon), and let $R$ be its
exterior obtained by deleting the six cycle vertices. Using signed port-count invariants
and simultaneous Kempe-switch closure across all $701$ six-port boundary types, we prove
that if $R$ is bipartite then $G$ admits a Berge--Fulkerson cover, and when each
bipartite class receives three ports then $G$ is $3$-edge-colorable. Moreover, if $G$ is
a smallest counterexample to the Berge--Fulkerson conjecture with six distinct exterior
terminals, then $R$ cannot become bipartite upon deleting any single internal edge or any
two incident internal edges, nor upon deleting two internal edges in one class when the
other class holds two opposite ports around the hexagon.
\end{abstract}

\section{Introduction}

Let $G$ be a cubic graph containing an induced $6$-cycle $C_6$ with vertices labeled
$0,1,2,3,4,5$ in cyclic order. Deleting the vertices of $C_6$ yields the \emph{exterior}
$R$ equipped with six corresponding semiedges (ports). A \emph{Berge--Fulkerson coloring}
assigns a $2$-element subset of six colors $\{0,\dots,5\}$ to each edge and semiedge so
that the three pairs incident with every vertex partition the six colors. In a closed
graph, the edges carrying each color form a perfect matching and every edge belongs to
two of the six matchings.

\begin{theorem}\label{thm:hexagon}
Let $G$ be a cubic graph with an induced $6$-cycle and exterior $R$.
\begin{enumerate}
  \item If $R$ is bipartite, then $G$ admits a Berge--Fulkerson cover; if in addition
        each bipartite class receives three ports, then $G$ is $3$-edge-colorable.
  \item Suppose $G$ is a smallest counterexample to the Berge--Fulkerson conjecture and
        the six exterior terminal vertices in $R$ are distinct. Then no $2$-partition
        $(+,-)$ of $V(R)$ can have:
        \begin{itemize}
          \item exactly one internal edge with both endpoints in the same class;
          \item exactly two internal edges within the classes sharing a common vertex; or
          \item exactly two internal edges within the $+$ class, none within the $-$
                class, and exactly two ports in the $-$ class at opposite positions
                around the hexagon.
        \end{itemize}
\end{enumerate}
\end{theorem}

\section{Boundary Types and Simultaneous Kempe Switches}

Because $3|V(R)| = 2|E(R)| + 6$, the exterior $R$ has even order. For any
Berge--Fulkerson coloring of $R$, each color $c$ covers $V(R)$ once by internal edges and
semiedges, so the subset of ports using color~$c$ has even cardinality. Recording the six
even port subsets and sorting them to forget color names produces the \emph{boundary type}
of the coloring. Every port appears in two of the six subsets; exact enumeration gives
exactly $701$ six-port boundary types. Colorings of $C_6$ and $R$ glue to a
Berge--Fulkerson cover of $G$ if and only if they realize a common boundary type.

On the hexagon $C_6$, a single color uses a matching of internal cycle edges together
with the semiedges at unmatched cycle vertices. Enumerating all six-element multisets of
such matchings that cover every cycle edge twice (which automatically gives every
semiedge multiplicity~$2$) shows that $C_6$ realizes exactly $92$ boundary types,
denoted $\Pi(C_6)$ (reproduced by \texttt{six\_caps.py}).

Fix a Berge--Fulkerson coloring of $R$ and two colors $a, b$. The subgraph of internal
edges and semiedges carrying exactly one of $a, b$ has degree $0$ or $2$ at every vertex
of $R$, so it decomposes into cycles and open paths pairing up the boundary ports that
carry exactly one of $a, b$. Swapping $a$ and $b$ along any subset of these open paths
yields another valid coloring of $R$. Moreover, for any partition of the six colors into
three disjoint pairs $\{a_1,b_1\}, \{a_2,b_2\}, \{a_3,b_3\}$, swaps on different pairs
act on disjoint color bits and therefore commute: even when paths for different pairs
share an edge, choosing an independent subset of paths for each pair and swapping all
three simultaneously produces a valid coloring of $R$.

This gives a sound closure operator on candidate supersets of the exterior profile
$\Pi(R)$, extending the five-cut switch method of~\cite{MMT}:

\begin{lemma}\label{lem:closure}
Let $A$ be a set of six-port boundary types such that $\Pi(R) \subseteq A$. Define the
simultaneous-switch pruning step by retaining a type $\tau \in A$ if and only if, for
every partition of the six colors into three disjoint pairs, there exists a choice of
pairings of the exclusive boundary ports for each color pair whose simultaneous
path-switch outcomes all lie in $A$. Let $\mathrm{cl}(A)$ be the fixed point obtained by
iterating this pruning step. Then $\Pi(R) \subseteq \mathrm{cl}(A)$.
\end{lemma}

\begin{proof}
Every actual type $\tau \in \Pi(R)$ is witnessed by a coloring of $R$, whose actual
bichromatic open-path pairings supply a valid joint pairing for every partition into
three color pairs with all simultaneous switch outcomes in $\Pi(R) \subseteq A$. By
induction on the iteration step, no element of $\Pi(R)$ is ever removed.
\end{proof}

Note that within any single color pair, switching all open chains merely exchanges the
names of its two colors without altering the unordered type (even under simultaneous
switches on the other two disjoint pairs), so it suffices to check subsets omitting one
chosen chain per pair.

In the unrestricted case (all $701$ types minus the $92$ hexagon types), one pruning step
reduces $609$ types to a fixed point of $561$ types. That fixed point still intersects
every orientation of eight tested smaller six-port caps (three wires, two junctions, an
edge and a wire, a four-vertex star, a four-vertex path, a triangle and junction, a
double edge and two wires, and a square and wire; logged in
\texttt{six-cap-closure.jsonl}), so simultaneous switches and these small caps alone do
not resolve unrestricted hexagons.

\section{Signed Port Counts and Bipartite Exteriors}

Partition $V(R)$ into two classes $+$ and $-$, and assume that every internal edge of $R$
crosses between $+$ and $-$ except for $t$ exceptional internal edges lying within the
$+$ class (and none within the $-$ class). Write $p$ for the number of ports attached to
vertices in $+$ (counted with multiplicity if terminals coincide) and set
$d = |+| - |-|$. Counting crossing edges from both sides gives
$3|+| - p - 2t = 3|-| - (6 - p)$, whence
\begin{equation}\label{eq:deg-count}
  3d = 2p - 6 + 2t.
\end{equation}
For any color $c$ in a Berge--Fulkerson coloring of $R$, let $m_c$ denote the number of
the $t$ exceptional edges carrying color~$c$. Counting crossing edges of color~$c$ from
both sides gives $|+| - \#(+\text{ ports using } c) - 2m_c = |-| - \#(-\text{ ports using } c)$,
and therefore
\begin{equation}\label{eq:signed-count}
  \#(+\text{ ports using } c) - \#(-\text{ ports using } c) = d - 2m_c.
\end{equation}
Across the six colors, the multiset $\{m_c\}_{c=0}^{5}$ is constrained as follows:
\begin{itemize}
  \item $t = 0$: $(0,0,0,0,0,0)$;
  \item $t = 1$: $(0,0,0,0,1,1)$;
  \item $t = 2$ with two incident exceptional edges: their labels are disjoint at the
        shared vertex, forcing $(0,0,1,1,1,1)$;
  \item $t = 2$ with two disjoint exceptional edges: letting their weight-$2$ labels
        overlap in $0$, $1$, or $2$ colors gives $(0,0,1,1,1,1)$, $(0,0,0,1,1,2)$, or
        $(0,0,0,0,2,2)$.
\end{itemize}

\begin{proof}[Proof of Theorem~\ref{thm:hexagon}(1)]
Suppose $R$ is bipartite ($t = 0$). By~\eqref{eq:deg-count}, $3d = 2p - 6$, so
$p \in \{0, 3, 6\}$. Up to dihedral symmetry of the hexagon and exchanging $+$ and $-$,
the set of $+$ ports has four representatives: $\varnothing$, $\{0,1,2\}$, $\{0,1,3\}$,
and $\{0,2,4\}$.

First, $R$ admits at least one $3$-edge-coloring (and hence, by doubling each color, at
least one Berge--Fulkerson coloring): since both $p$ and $6 - p$ are multiples of~$3$,
group same-class ports into triples and attach each triple to a new vertex placed in the
opposite class. The resulting graph is a $3$-regular bipartite multigraph (with
$|+| + (6-p)/3 = |-| + p/3$ vertices per side), which decomposes into three perfect
matchings by Hall's/K\H{o}nig's theorem. Restricting to $R$ gives a $3$-edge-coloring of
$R$.

For Berge--Fulkerson covers, let $E_0$ be the subset of the $701$ six-port types
satisfying~\eqref{eq:signed-count} with $m_c = 0$ for all $c$. Exact computation in
\texttt{six\_caps.py} shows that $\mathrm{cl}(E_0 \setminus \Pi(C_6)) = \varnothing$ in
all four cases (iteration sizes: $119, 94, 20, 0$ for $\varnothing$; $153, 122, 38, 0$ for
$\{0,1,2\}$; $144, 55, 0$ for $\{0,1,3\}$; and $77, 0$ for $\{0,2,4\}$). Since
$\Pi(R) \neq \varnothing$, Lemma~\ref{lem:closure} forces
$\Pi(R) \cap \Pi(C_6) \neq \varnothing$, so $G$ has a Berge--Fulkerson cover.

When $p = 3$ (so $d = 0$), apply the same method directly to $3$-edge-colorings without
doubling. A $3$-color boundary type is an unordered partition of the six ports into three
even color supports, each containing equally many $+$ and $-$ ports. For each of the
three $p = 3$ representatives, there are $16$ eligible $3$-color types. Computing the
$3$-color boundary types realized by proper $3$-edge-colorings of $C_6$ and applying the
single-pair Kempe-switch closure removes all types avoiding the hexagon: the surviving
counts drop as $11, 8, 2, 0$ for $\{0,1,2\}$; $10, 3, 0$ for $\{0,1,3\}$; and $5, 0$ for
$\{0,2,4\}$. Since $R$ admits at least one $3$-edge-coloring, its $3$-color profile meets
that of $C_6$, so $G$ is $3$-edge-colorable.
\end{proof}

\section{Near-Bipartite Exteriors and Bridgeless Wire Completions}

We now prove Theorem~\ref{thm:hexagon}(2). Assume $G$ is a smallest counterexample
(minimum order among simple bridgeless cubic graphs with no Berge--Fulkerson cover) with
an induced $6$-cycle $C_6$ and six distinct exterior terminals in $R$. By standard
small-cut reductions~\cite{MMT}, $G$ is cyclically $5$-edge-connected.

\begin{lemma}\label{lem:wire-bridgeless}
The exterior $R$ is connected and bridgeless. Moreover, closing $R$ by any matching of
three wire edges pairing up its six distinct terminals produces a smaller bridgeless
cubic multigraph that admits a Berge--Fulkerson cover.
\end{lemma}

\begin{proof}
Because the six exterior terminals are distinct, every vertex of $R$ has internal degree
at least~$2$, so every connected component of $R$ contains a cycle. Two components would
each require at least $5$ ports (since each complement includes $C_6$), impossible with
six ports total; hence $R$ is connected.

Suppose $R$ had an internal bridge $e = uv$ separating $R$ into components $R_1$ and
$R_2$. In each $R_i$, every vertex other than its endpoint of $e$ has internal degree at
least~$2$, and that endpoint has degree at least~$1$ in $R_i$. A tree on at least two
vertices has at least two leaves, and a $1$-vertex component has degree~$0$, so neither
$R_1$ nor $R_2$ can be a tree; both contain cycles. Since each complement in $G$ contains
$C_6$, cyclic $5$-edge-connectivity forces each $R_i$ to carry at least $4$ ports in
addition to $e$, requiring at least $8$ ports total, a contradiction. Thus $R$ is
bridgeless.

Adding three wires matching the six distinct terminals gives a cubic multigraph $G'$ on
$|V(G)| - 6$ vertices with no loops. Every edge of $R$ lies on a cycle in $R$, and each
wire closes to a cycle via a path in the connected graph $R$, so $G'$ is bridgeless.

It remains to verify by strong induction on $|V(H)|$ that every bridgeless cubic
multigraph $H$ with $|V(H)| < |V(G)|$ admits a Berge--Fulkerson cover. Work component by
component; no component contains a loop (whose third incident edge would be a bridge). If
$H$ is simple, minimality of $G$ supplies a cover. If $H$ consists of a triple edge,
partition the six colors into three pairs on its three edges. Otherwise $H$ contains a
double edge $xy$ with external edges $xx'$ and $yy'$. We must have $x' \neq y'$, or else
the third edge at $x' = y'$ would be a bridge. Deleting $x, y$ and adding a single edge
$e^* = x'y'$ produces a strictly smaller cubic multigraph $H'$ that remains bridgeless
(any cycle in $H$ through $xx'$ traverses one edge of $xy$ and exits via $yy'$, yielding
a cycle through $e^*$ in $H'$, and all other cycles persist). By induction, $H'$ admits a
cover; assigning the label of $e^*$ to both $xx'$ and $yy'$ and splitting the remaining
four colors across the two parallel edges $xy$ lifts the cover to $H$. Applying this to
$H = G'$ completes the proof.
\end{proof}

\begin{proof}[Proof of Theorem~\ref{thm:hexagon}(2)]
Consider a partition $(+,-)$ of $V(R)$ with $t \in \{1,2\}$ exceptional edges inside $+$.
By~\eqref{eq:deg-count}:
\begin{itemize}
  \item For $t = 1$, $3d = 2p - 4$, so $p \in \{2, 5\}$. Up to dihedral symmetry of $C_6$,
        the $+$ port sets have representatives $\{0,1\}$, $\{0,2\}$, $\{0,3\}$, and
        $\{0,1,2,3,4\}$.
  \item For $t = 2$, $3d = 2p - 2$, so $p \in \{1, 4\}$. Up to dihedral symmetry, the $+$
        port sets have representatives $\{0\}$, $\{0,1,2,3\}$, $\{0,1,2,4\}$, and
        $\{0,1,3,4\}$.
\end{itemize}
For $t = 1$ and for two incident exceptional edges, let $E$ be the set of types allowed
by~\eqref{eq:signed-count} with multiplicities $(0,0,0,0,1,1)$ and $(0,0,1,1,1,1)$,
respectively. For two (possibly disjoint) exceptional edges in $+$ with opposite $-$
ports (represented by $+$ ports $\{0,1,3,4\}$), let $E$ be the union of the type sets
allowed by~\eqref{eq:signed-count} over all three label overlaps ($0$, $1$, or $2$ shared
colors). Since a Kempe switch may change the overlap between two disjoint exceptional
edges, taking the union before closure ensures $\Pi(R) \subseteq E$.

If $G$ has no Berge--Fulkerson cover, then $\Pi(R) \subseteq E \setminus \Pi(C_6)$, and
hence $\Pi(R) \subseteq \mathrm{cl}(E \setminus \Pi(C_6))$ by Lemma~\ref{lem:closure}.
For each row of Table~\ref{tab:wires}, \texttt{six\_caps.py} computes the exact fixed
point $\mathrm{cl}(E \setminus \Pi(C_6))$ and enumerates the profile of the listed
three-wire completion (all types whose port labels agree across each wire). In every row,
the fixed point is disjoint from the wire completion's profile. However, by
Lemma~\ref{lem:wire-bridgeless}, closing $R$ with that wire completion yields a smaller
bridgeless cubic multigraph admitting a Berge--Fulkerson cover, whose restriction to $R$
realizes a type in the wire completion's profile, contradicting
$\Pi(R) \subseteq \mathrm{cl}(E \setminus \Pi(C_6))$.
\end{proof}

\begin{table}[ht]
\centering
\caption{Simultaneous-switch closure sizes and separating wire completions for near-bipartite exteriors.}
\label{tab:wires}
\begin{tabular}{@{}llrrl@{}}
\toprule
Exceptional edges & Positive ports & Eligible & Surviving types & Wire completion \\
\midrule
One & $0,1$         & $522$ & $368$ & $01,23,45$ \\
One & $0,2$         & $522$ & $0$   & $01,23,45$ \\
One & $0,3$         & $522$ & $248$ & $03,12,45$ \\
One & $0,1,2,3,4$   & $425$ & $222$ & $01,25,34$ \\
Two sharing a vertex & $0$         & $425$ & $222$ & $03,12,45$ \\
Two sharing a vertex & $0,1,2,3$   & $522$ & $368$ & $01,23,45$ \\
Two sharing a vertex & $0,1,2,4$   & $522$ & $0$   & $01,23,45$ \\
Two sharing a vertex & $0,1,3,4$   & $522$ & $248$ & $01,25,34$ \\
Two in one class     & $0,1,3,4$   & $649$ & $286$ & $01,25,34$ \\
\bottomrule
\end{tabular}
\end{table}

\section{Reproduction}

Run \texttt{python3 six\_caps.py} from the repository root to reproduce all boundary-type
counts, simultaneous-switch closure sequences, and separating wire completions reported
above (archived in \texttt{six-cap-closure.jsonl}).

\begin{thebibliography}{9}
\bibitem{MMT}
E.~M\'{a}\v{c}ajov\'{a}, G.~Mazzuoccolo, and G.~Tabarelli,
Ars Math.\ Contemp., \href{https://doi.org/10.26493/1855-3974.3409.c13}{doi:10.26493/1855-3974.3409.c13}.
\end{thebibliography}

\end{document}
